\section{Design-space exploration}
\label{sec:dse}

\subsection{The DSE driver}
A \emph{campaign} (\file{configs/dse/*.json}, \cref{lst:campaign}) names a base configuration, a build grid of
dotted configuration keys with value lists, a runtime grid (matrix sizes as lists or explicit triples, repetitions,
warm-up, verification mode) and the stages to run. \code{scripts/gemm\_dse.py} (\cref{fig:dse}):
\begin{enumerate}
  \item expands the full cross product of the build grid (values can be objects, e.g.\ whole memory maps) and names
        each point from its parameters;
  \item validates every point and records invalid ones with the reason, without running any tool;
  \item \textbf{hls} stage: runs HLS for all points in parallel, as many jobs as \(\min(\text{cores}/4,
        \text{free GB}/3)\) allows, and records the estimates (clock, LUT, FF, DSP, BRAM);
  \item \textbf{build} stage: runs LiteX and Vivado for one point at a time; a point whose routed design misses timing
        on \emph{any} clock is recorded as \code{timing\_failed} and is never loaded onto the card;
  \item \textbf{run} stage: for each point with timing met, programs the card and brings it up (JTAG, rescan,
        driver, DDR calibration, host tools), then runs every size of the runtime grid.
\end{enumerate}
Results go to \file{results/<campaign>/}: \code{builds.jsonl} (one row per point and stage), \code{runs.jsonl} (one row
per point, size and repetition, with the full configuration, the hashes and the git state), regenerated CSV files with
flattened columns, and \code{invocations.jsonl} (every invocation of the campaign). \code{--resume} skips stages and
sizes already recorded as \code{ok}, and does not even reprogram the card for a point with nothing left to run; it was
tested by deleting rows from a finished campaign and resuming on the card.

\lstinputlisting[language={}, caption={Campaign \file{configs/dse/pe32\_ta\_225\_presets.json}: two Vivado presets
  at 225\,MHz, three sizes.}, label={lst:campaign}, basicstyle=\ttfamily\scriptsize]{../../configs/dse/pe32_ta_225_presets.json}

\begin{figure}[t]
  \centering
  \resizebox{0.95\linewidth}{!}{% The DSE driver (scripts/gemm_dse.py) and the outer design loop around it.
\begin{tikzpicture}[x=1cm,y=1cm, every node/.append style={font=\sffamily\footnotesize}]
  \tikzset{st/.style={flw, minimum width=29mm, minimum height=12mm, align=center}}
  \node[art, text width=30mm, minimum height=13mm] (camp) at (0,0)
       {\textbf{campaign JSON}\\ base config, \code{build\_grid}, \code{runtime\_grid}, stages};
  \node[st] (exp) at (3.9,0)  {expand\\[-1pt]\scriptsize cross product of\\[-1pt]\scriptsize dotted keys};
  \node[st] (val) at (7.6,0)  {validate\\[-1pt]\scriptsize gemm\_hls rules;\\[-1pt]\scriptsize invalid points recorded};
  \node[st] (hls) at (11.3,0) {stage \textbf{hls}\\[-1pt]\scriptsize parallel: min(cores/4,\\[-1pt]\scriptsize free GB/3) jobs};
  \node[st] (bld) at (11.3,-2.2) {stage \textbf{build}\\[-1pt]\scriptsize Vivado, one at a time;\\[-1pt]\scriptsize \code{ok} / \code{timing\_failed}};
  \node[st] (run) at (7.6,-2.2)  {stage \textbf{run}\\[-1pt]\scriptsize per build: bring-up,\\[-1pt]\scriptsize sizes \x repetitions};
  \node[art, text width=33mm, minimum height=13mm] (res) at (3.9,-2.2)
       {\file{results/<campaign>/}\\ \code{builds.jsonl}, \code{runs.jsonl} + CSV, \code{invocations.jsonl}};
  \draw[link] (camp) -- (exp); \draw[link] (exp) -- (val); \draw[link] (val) -- (hls);
  \draw[link] (hls) -- (bld); \draw[link] (bld) -- node[lbl, above] {timing met} (run); \draw[link] (run) -- (res);
  \node[dec] (an) at (0,-2.2) {analyse};
  \draw[link] (res) -- (an);
  \draw[link, cRed!80!black, very thick] (an.north) -- node[lbl, left, text=cRed!80!black, align=right]
       {next\\campaign} (camp.south);
  % reuse and resume annotations
  \node[note, anchor=north, text width=58mm, align=left] at (9.45,-3.35)
       {\(\bullet\) point names come from the parameters (\code{f32\_pe32x8\_t1024x512\_225mhz\_tA\_aggressive}),
        so equal points share one build directory across campaigns\\
        \(\bullet\) unchanged steps are skipped from their stamps; \code{--resume} skips recorded
        stages and sizes (and the bring-up of a point with nothing left)\\
        \(\bullet\) a build that misses timing anywhere is never loaded onto the card};
  \node[note, anchor=north, text width=40mm, align=left] at (1.9,-3.35)
       {The analysis is done by the engineer (here with an AI assistant): read the counters and the timing paths,
        pick the next knob.};
\end{tikzpicture}
}
  \caption{The DSE driver and the design loop around it. Each campaign's analysis picks the next campaign.}
  \label{fig:dse}
\end{figure}

\subsection{The design loop}
The driver runs a campaign without supervision; choosing the next campaign is the engineer's job, guided by the
counters and the timing reports. \Cref{fig:journey} shows how the exploration actually went, and
\cref{tab:campaigns} lists every campaign with its size and duration, taken from the result records. Builds are shared
between campaigns: the same point appears in several of them but is built once.

\begin{figure}[t]
  \centering
  \resizebox{\linewidth}{!}{% The design loop as it happened: campaign -> finding -> next knob, with the best result after each step.
\begin{tikzpicture}[x=1cm,y=1cm, every node/.append style={font=\sffamily\footnotesize}]
  \tikzset{cp/.style={flw, minimum width=31mm, minimum height=10mm, align=center, text width=29mm},
           fd/.style={note, text width=31mm, align=left, anchor=north},
           gb/.style={draw=cGreen!60!black, fill=cGreen!12, rounded corners=6pt, font=\sffamily\scriptsize\bfseries,
                      inner sep=2pt, anchor=south, yshift=0.8mm}}
  \node[cp] (c1) at (0,0)    {\code{pe\_sweep}\\[-1pt]\scriptsize \(P_N\) = 8, 16, 32 (tile 256)};
  \node[cp] (c2) at (3.9,0)  {\code{transposed\_a}\\[-1pt]\scriptsize host sends A\(^\mathsf{T}\)};
  \node[cp] (c3) at (7.8,0)  {\code{tiles\_pe32\_ta}\\[-1pt]\scriptsize \code{\_1024}: tiles up to 1024\x512};
  \node[cp] (c4) at (11.7,0) {\code{clock\_sweep}, presets\\[-1pt]\scriptsize PE 8\x8: 150--250\,MHz};
  \node[cp] (c5) at (15.6,0) {\code{clock\_pe32\_ta},\\[-1pt]\code{pe32\_ta\_225\_presets}};
  \foreach \a/\b in {c1/c2, c2/c3, c3/c4, c4/c5} { \draw[link, thick] (\a) -- (\b); }
  \node[gb] at (c1.north) {86.5 GFLOP/s};
  \node[gb] at (c2.north) {87.5};
  \node[gb] at (c3.north) {95.1 \(\rightarrow\) 98.6};
  \node[gb] at (c4.north) {PE 8\x8: 29.0 at 250\,MHz};
  \node[gb] at (c5.north) {110.9 at 225\,MHz};
  \node[fd] at ([yshift=-1mm]c1.south) {\textbf{Finding:} PE 32\x8 reaches 84\,\% only; A waits 93\,\% of cycles: A is
        read in single 64-byte strided beats.};
  \node[fd] at ([yshift=-1mm]c2.south) {\textbf{Finding:} A now in 16-beat bursts; A and B wait equally (16\,\%).
        The per-tile reload of A/B is next.};
  \node[fd] at ([yshift=-1mm]c3.south) {\textbf{Finding:} larger tiles halve the waits; the N tile matters most. 1024\x1024
        does not fit in BRAM.};
  \node[fd] at ([yshift=-1mm]c4.south) {\textbf{Finding:} PE 8\x8 closes 250\,MHz despite a 4.24\,ns HLS estimate; the only
        failure was a PCIe path; presets change closure only.};
  \node[fd] at ([yshift=-1mm]c5.south) {\textbf{Finding:} 225\,MHz misses by 41\,ps with the default preset (routing);
        \code{explore}/\code{aggressive} close it. 250\,MHz: \(-0.78\)\,ns.};
\end{tikzpicture}
}
  \caption{The exploration as it happened. Each campaign's finding decided the next knob; badges show the best
    result of the campaign (kernel GFLOP/s).}
  \label{fig:journey}
\end{figure}

\begin{table}[t]
\centering\small
\caption{The float32 campaigns, 23--24 November 2026 (\file{docs/report/data/campaigns.csv}). Points: design points;
  builds: new Vivado runs (reused builds not counted); runs: verified result rows. The integer and fixed-point
  campaigns are in \cref{tab:fmt-campaigns}.}
\label{tab:campaigns}
\begin{tabular}{@{}lllrrrl@{}}
\toprule
Campaign & Time & Stages & Points & Builds & Runs & Question \\
\midrule
\code{hls\_pe\_tile} & 17:08--17:10 & hls & 12 & -- & -- & HLS estimates over PEs and tiles \\
\code{pe\_sweep} & 17:11--18:05 & all & 3 & 3 & 24 & how does throughput scale with \(P_N\)? \\
\code{transposed\_a} & 18:08--18:50 & all & 2 & 2 & 16 & does reading A in bursts help? \\
\code{tiles\_pe32\_ta} & 18:51--20:10 & all & 4 & 3 & 24 & which memory tile? \\
\code{pe32\_ta\_next\_hls} & 20:11--20:14 & hls & 8 & -- & -- & do larger tiles fit? 225\,MHz? \\
\code{tiles\_pe32\_ta\_1024} & 20:15--20:43 & all & 2 & 1 & 12 & tile 1024\x512 \\
\code{memory\_maps} & 20:47--21:12 & all & 2 & 2 & 16 & other channel assignments \\
\code{size\_sweep} & 21:13--21:14 & build, run & 1 & -- & 90 & 45 matrix shapes on the best build \\
\code{clock\_sweep} & 21:15--21:54 & all & 3 & 3 & 4 & PE 8\x8 at 150/200/250\,MHz \\
\code{vivado\_preset\_sweep} & 21:54--22:21 & build, run & 3 & 2 & 2 & presets at 200\,MHz \\
\code{clock\_pe32\_ta} & 22:26--23:30 & all & 3 & 2 & 6 & best config at 200/225/250\,MHz \\
\code{pe32\_ta\_225\_presets} & 23:31--00:41 & build, run & 2 & 2 & 12 & presets at 225\,MHz \\
\midrule
\multicolumn{3}{@{}l}{total} & & 20 & 206 & \\
\bottomrule
\end{tabular}
\end{table}
